% 符号对照表。规范里这一节可以放多张表，各有各的列数，所以这里不做成固定的两列
% 环境，就是普通表格，写几张、每张几列自己定。
%
% 用 longtable 而不是 table，条目多了会自动跨页；table 是浮动体，装不下就整个溢出。
% 长表格的双语标题只能用表格的前两行写，所以走 \longbionenumcaption，七个参数依次是
% 中文短标题、中文标题、英文表名、英文短标题、英文标题、中英标题间距、标题与表的间距。
% 只要中文标题的话，把那一行换成 \caption{标题}\\ 就够。

\begin{denotation}
\ltfontsize{\wuhao[1.667]}
\wuhao[1.667]\begin{longtable}{cccc}%
\longbionenumcaption{}{{\wuhao 国际单位制中具有专门名称的导出单位}\label{table:si}}{Table$\!$}{}{{\wuhao Derived units with special names in SI}}{-0.5em}{3.15bp}\\
\toprule 量的名称 & 单位名称 & 单位符号 & 其它表示实例\\ \midrule
\endfirsthead
\multicolumn{4}{r}{表~\thetable（续表）}\vspace{0.5em}\\
\toprule 量的名称 & 单位名称 & 单位符号 & 其它表示实例\\ \midrule
\endhead
\midrule[0.5pt]
\endfoot
\bottomrule
\endlastfoot
频率 & 赫[兹] & Hz & s$^{-1}$\\
力 & 牛[顿] & N & kg$\cdot$m/s$^2$\\
压力 & 帕[斯卡] & Pa & N/m$^2$\\
能量 & 焦[耳] & J & N$\cdot$m\\
功率 & 瓦[特] & W & J/s\\
\end{longtable}
\end{denotation}
